\documentclass{amsart}
\usepackage{graphicx,hyperref}
\newtheorem{theorem}{Theorem}[section]
\newtheorem{lemma}[theorem]{Lemma}
\theoremstyle{definition}
\newtheorem{definition}[theorem]{Definition}
\theoremstyle{remark}
\newtheorem{remark}[theorem]{Remark}
\numberwithin{equation}{section}
\begin{document}
\title[Synthetic audit]{Synthetic AMS compatibility audit}
\author{Alpha Author}
\address{Department Alpha, Synthetic Institute One}
\curraddr{Current Department Alpha, Synthetic Institute One}
\email{alpha@example.invalid}
\thanks{Synthetic support statement for Alpha.}
\author{Beta Author}
\address{Department Beta, Synthetic Institute Two}
\email{beta@example.invalid}
\thanks{Synthetic support statement for Beta.}
\subjclass[2020]{Primary 05C50; Secondary 15A18}
\keywords{bounded graph, spectral inequality}
\date{}
\dedicatory{Dedicated to synthetic examples}
\begin{abstract}
We prove a synthetic bound for a two-vertex graph. This abstract is original audit text.
\end{abstract}
\maketitle
\section{The bound}\label{sec:bound}
\begin{definition}[Graph]\label{def:graph}
Let $G$ have two vertices and one edge.
\end{definition}
\begin{theorem}[Energy]\label{thm:energy}
For $x,y\in\mathbb{R}$, $(x-y)^2\geq 0$.
\end{theorem}
\begin{proof}
Expand the square in equation~\eqref{eq:energy}.
\end{proof}
\begin{align}
E(x,y)&=(x-y)^2\label{eq:energy}\\
&=x^2-2xy+y^2.\label{eq:expand}
\end{align}
\begin{lemma}\label{lem:zero}
If $x=y$, the energy vanishes.
\end{lemma}
\begin{remark}
This is a synthetic example, cited as~\cite{synthetic}.
\end{remark}
\begin{figure}[ht]
\centering\includegraphics[width=25mm]{audit.png}
\caption{Synthetic two-color audit image}\label{fig:image}
\end{figure}
\begin{table}[ht]
\caption{Synthetic values}\label{tab:values}
\centering\begin{tabular}{lc}\hline Case&Energy\\\hline Equal&0\\Distinct&1\\\hline\end{tabular}
\end{table}
By Theorem~\ref{thm:energy} and Lemma~\ref{lem:zero}, see Figure~\ref{fig:image} and Table~\ref{tab:values}.
\appendix
\section{Additional identity}\label{sec:appendix}
\begin{equation}\label{eq:appendix}1+1=2.\end{equation}
\bibliographystyle{amsplain}
\bibliography{references}
\end{document}
